%% source: 05_experiments (11.3kw) -> target ~2kw

\section{Results}\label{sec:results}

\subsection{Experimental setup and corpora}\label{sec:results:setup}
We evaluate the pipeline across three tiers of breadth.
A deep benchmarking corpus of 13 MapLibre styles is rendered end-to-end against the Germany OpenMapTiles (OMT) tileset (zoom levels 0 to 14, \qty{256.5}{\kilo tiles}) under \num{18} stratified scenarios, each a pairing of a geographic location with a camera animation, selected so that every location appears at least once and dense and sparse, Latin and non-Latin cases are all represented.
A second tier of \num{11} styles whose referenced tile data is fully available is evaluated at the tile-shaving size level on the same tileset, and a third generalization corpus of \num{13} additional unseen styles drawn from the Maputnik gallery is evaluated with static analysis only.
The measurement harness is a Puppeteer-driven headless browser that loads MapLibre GL JS, replays each animation, and collects load time, frames per second (FPS), frame-time percentiles, and time to first tile through the browser \texttt{Performance} API with real GPU rendering (ANGLE/Vulkan).
Each style and scenario cell is run \num{15} times and reported as the median.
The pipeline itself is measured through a 19-step cumulative ablation in which each pass is added on top of all preceding passes, plus one plain-tile-shaving comparison configuration.
Three styles using legacy \texttt{text-field} syntax that an upstream MapLibre migrator miscompiles are capped at the structural stage and excluded from the end-to-end pool of 10 styles.

\subsection{Style-level optimization}\label{sec:results:style}
The eight expression-level passes barely moved end-to-end load time on their own, reducing the per-style median by \qty{2.0}{\percent} (range \qtyrange{-58.8}{+12.3}{\percent} across styles), an effect not consistently distinguishable from per-style measurement noise.
Their effect on serialized style size was clearer: default stripping was the single largest marginal contributor at \qty{-4.3}{\percent} of raw style size (\qty{-2.9}{\percent} under gzip), followed by stats-driven folding (\qty{-1.6}{\percent}) and expression simplification (\qty{-1.4}{\percent}).
Raw style size fell by up to \qty{20.3}{\percent} on Fiord (\num{22.2} to \qty{17.7}{\kilo\byte}) and \qty{7.0}{\percent} on Liberty, narrowing to \qty{8.3}{\percent} and \qty{5.9}{\percent} respectively under gzip as the compressor absorbed part of the raw savings.
The size reduction translated directly into a shorter style-parse step even where end-to-end load time was unchanged.

The structural passes concentrated almost all of their gains in dead-layer elimination and layer merging, with the remaining passes negligible in isolation.
The magnitude was style-dependent: Fiord's 48-layer style shed two layers (\qty{-4.2}{\percent}, \qty{-14.1}{\percent} gzip transfer, \qty{-36.8}{\percent} expression-tree nodes from \num{1610} to \num{1018}), while Liberty's 111-layer style shed four through layer merging (\qty{-3.6}{\percent}, \qty{-4.9}{\percent} gzip, nodes from \num{3537} to \num{3167}).
Layer merging can increase raw byte count, because a merged layer replaces several short filters with wordier \texttt{match}-keyed expressions, yet still reduces per-frame draw calls and frame time.
Gzip absorbs the raw increase, leaving compressed size essentially unchanged.
Symbol-heavy styles benefit less because symbol layers are excluded from merging on collision-soundness grounds.

\subsection{Tile shaving and columnar re-encoding}\label{sec:results:data}
Tile shaving was the dominant data-level lever.
Pooled across the \num{10} deep-corpus styles, shaving alone cut the median load time from \qty{390}{\milli\second} to \qty{135}{\milli\second} (\qty{-65.5}{\percent}) and more than doubled median FPS (\qty{+113.2}{\percent}), whereas style-only optimization contributed a modest \qty{-4.0}{\percent} (Figure~\ref{fig:box}).
Across the \num{11} styles with available tile data, style-driven on-wire reductions on the Germany tileset ranged from \qty{19.8}{\percent} (Bright) to \qty{46.3}{\percent} (Toner) with a median of \qty{29.0}{\percent} (Figure~\ref{fig:tile_shave_per_style}), the spread reflecting how many source-layers and property columns each style references.

\begin{figure}[tb]
    \centering
    \includegraphics[width=\linewidth]{tile_shave_per_style.pdf}
    \caption{On-wire data reduction from style-driven MVT shaving across the \num{11} styles with fully available tile data, applied to the same \qty{3.07}{\giga\byte} Germany tileset.
    Styles that reference fewer source-layers and property columns achieve larger reductions.}
    \label{fig:tile_shave_per_style}
\end{figure}

The columnar MapLibre Tile (MLT) format alone was not the win.
Table~\ref{tab:encoder} reports the total encoded size of the Germany OMT tileset under three encoders.
The reference MLT encoder drifted to \qty{3.54}{\giga\byte} under gzip, worse than gzip-compressed Mapbox Vector Tiles (MVT) at \qty{3.08}{\giga\byte}, because general-purpose compressors already exploit most row-level redundancy that a fixed-plan columnar encoder cannot match.
Our encoder, which adds sort-strategy competition, per-stream profiling, and shared-dictionary clustering, reached \qty{3.06}{\giga\byte} uncompressed, already below MVT+gzip, and \qty{2.70}{\giga\byte} under gzip: \qty{28.4}{\percent} smaller than the reference baseline uncompressed and \qty{12.3}{\percent} smaller than MVT+gzip.
The advantage held at \qty{13.0}{\percent} even against zstd-22, the strongest general-purpose compressor tested.
On the decode axis, Criterion micro-benchmarks of our native Rust MLT decoder showed \numrange{4.6}{8.3}$\times$ higher full-decode throughput than MVT+gzip across zoom levels 4, 7, and 13 (Table~\ref{tab:encoder}), corresponding to \numrange{96}{134} versus \numrange{15}{23}~MiB/s.
This decode advantage does not appear in the browser load-time numbers, which use the existing JavaScript MLT decoder rather than the Rust path.
Decomposing the on-wire savings by zoom level (Figure~\ref{fig:shaving}) separates the two effects: the gap between each format and its shaved counterpart is the payload removed by the pruning advisory, which dominates at the dense mid-zoom levels, while the residual gap between shaved MVT and shaved MLT is the additional contribution of columnar encoding.

\begin{table}[tb]
\centering
\caption{Encoded size and full-decode latency for the Germany OMT tileset (\qty{256.5}{\kilo tiles}).
The reference MLT column carries our backported encoding fixes, isolating the gain from automatic strategy selection.
Latency is Criterion full-decode time (100 samples per cell) of our Rust decoder versus MVT+gzip, gzip decompression included (uncompressed MVT: \qtylist{44.2;90.6;16.5}{\milli\second}).
The vs.\ MVT row gives size reduction for the size columns and decode speedup for the latency columns.}
\label{tab:encoder}
\setlength{\tabcolsep}{4.5pt}
\begin{tabular}{l rrrr rrr}
\toprule
& \multicolumn{4}{c}{Encoded size (GB)} & \multicolumn{3}{c}{Full-decode latency (ms)} \\
\cmidrule(lr){2-5} \cmidrule(lr){6-8}
Encoder & Plain & Gzip & Brotli & Zstd & Zoom 4 & Zoom 7 & Zoom 13 \\
\midrule
MVT & 4.24 & 3.08 & 3.13 & 3.16 & 47.6 & 103 & 18.7 \\
MLT (reference) & 4.26 & 3.54 & 3.57 & 3.60 & --- & --- & --- \\
MLT (this work) & \textbf{3.06} & \textbf{2.70} & \textbf{2.72} & \textbf{2.75} & \textbf{5.73} & \textbf{19.4} & \textbf{4.09} \\
\midrule
$\Delta$ vs.\ reference & \num{-28.4}\% & \num{-23.8}\% & \num{-23.7}\% & \num{-23.6}\% & --- & --- & --- \\
vs.\ MVT & \num{-28.2}\% & \num{-12.3}\% & \num{-13.1}\% & \num{-13.0}\% & $8.3\times$ & $5.3\times$ & $4.6\times$ \\
\bottomrule
\end{tabular}
\end{table}

\begin{figure}[tb]
    \centering
    \includegraphics[width=0.85\linewidth]{shaving_effectiveness_per_zoom}
    \caption{Total tile size per zoom level for unshaved and shaved variants of both MVT and our MLT encoder, using the Liberty advisory.
    The gap between each format and its shaved counterpart is the data removed by the pruning advisory, and the gap between shaved MVT and shaved MLT is the additional contribution of columnar encoding.}
    \label{fig:shaving}
\end{figure}

\subsection{Combined pipeline: additive or synergistic?}\label{sec:results:combined}
The central empirical question was whether style-level and data-level optimizations interact synergistically, and the answer depends on the level and the style.
On data volume the two are cleanly additive for every style: across the \num{10} styles, shaving the optimized style removed essentially the same tile bytes as shaving the unoptimized style, the optimized advisory adding only \qtyrange{0.0}{0.5}{\percent} of extra reduction (median \qty{0.08}{\percent}).
Style simplification does not narrow what the pruning advisory removes, because the advisory turns on which source-layers and property columns the style references at all, which simplification leaves intact.
On end-to-end load time the levels are additive in aggregate as well: pooled across the \num{10} styles, the combined reduction (\qty{66.4}{\percent}) approximately equals the sum of the style-only (\qty{4.0}{\percent}) and shaving-only (\qty{65.5}{\percent}) reductions, with an interaction term of \qty{-3.1}{\percent}.
Per scenario, however, small super-additive interactions are common rather than rare, and they vary sharply by style: combined load time beat the sum of the parts in \num{103} of \num{180} style-scenario pairs (\qty{57}{\percent}).
The effect is concentrated in a subset of styles (Bright \num{15} of \num{18} pairs, median interaction \qty{+7.1}{\percent}; OSM~Bright and OSM~Liberty \num{15} of \num{18}; Stadia~Outdoors \num{14}; Liberty \num{13}) and largely absent in others (Fiord \num{4} of \num{18}, Positron \num{5}, Dark~Matter \num{6}, Klokan-Basic \num{8}, Toner \num{8}).
This split does not track how much each style shaves: Toner shaves the most tile data yet stays additive, while OSM~Bright shaves the least yet synergizes the most.
The interactions are also small in magnitude (median \qty{+0.64}{\percent} across all pairs), and because the sub-additive styles offset the super-additive ones they do not lift the aggregate.
Since tile volume is unchanged by style optimization, this residual per-scenario synergy is a second-order rendering and decode timing effect rather than a data-volume effect.
We therefore read style-level passes as enablers that simplify the advisory, and for some styles modestly improve render timing, rather than as co-equal synergistic partners.
The full 19-step ablation confirms this picture (Figure~\ref{fig:box} and Figure~\ref{fig:waterfall}): load-time variance compresses sharply only when tile shaving fires, the FPS spread widens past the same step as some styles realize more headroom, and the wire-size waterfall shows almost all of the style-size win landing by the metadata-stripping step.
A breakdown of where the load-time budget goes across the ablation (Figure~\ref{fig:time_breakdown}) confirms this division of labor: the style-parse leg shrinks monotonically with each expression pass, while the first-tile and remaining-load legs are untouched by style-only steps and collapse only once tile shaving fires.
Composed end-to-end across the 10 deep-corpus styles, the full pipeline reduced load time by \qtyrange{63}{74}{\percent} (per-style median \qty{71}{\percent}) and improved FPS by \qtyrange{100}{200}{\percent} (per-style median \qty{152}{\percent}).
Browser heap memory fell \qty{35.7}{\percent} from \qty{86.0}{\mega\byte} to \qty{55.3}{\mega\byte}, with almost all of the reduction coming from the smaller in-memory tile payload after shaving (Figure~\ref{fig:memory_ablation}).

\begin{figure}[tb]
    \centering
    \includegraphics[width=\linewidth]{time_breakdown.pdf}
    \caption{Load time decomposed into style parse, first tile, first frame, and remaining load across the ablation steps.
    Style-parse time shrinks monotonically with each expression pass, while first-tile and remaining-load times are unchanged by style-only steps and decrease only when tile shaving is applied.}
    \label{fig:time_breakdown}
\end{figure}

\begin{figure}[tb]
    \centering
    \includegraphics[width=\linewidth]{memory_ablation.pdf}
    \caption{Mean browser heap usage across all scenarios and benchmark styles by ablation step.
    Style-level passes trim a few megabytes, but the bulk of the \qty{-35.7}{\percent} reduction comes from the data-level passes, where the in-memory tile payload itself shrinks after shaving.}
    \label{fig:memory_ablation}
\end{figure}
Bootstrap \qty{95}{\percent} confidence intervals confirmed the full-pipeline load-time reductions were tight (Fiord \qty{370.9}{\milli\second} to \qty{120.5}{\milli\second}, Liberty \qty{466.5}{\milli\second} to \qty{133.7}{\milli\second}, both $p < 10^{-45}$, Wilcoxon).

\begin{figure}[tb]
    \centering
    \begin{subfigure}[t]{\linewidth}
        \centering
        \includegraphics[width=\linewidth]{box_loadMs.pdf}
        \caption{Load time per ablation step.}
        \label{fig:box_loadMs}
    \end{subfigure}

    \vspace{0.5em}

    \begin{subfigure}[t]{\linewidth}
        \centering
        \includegraphics[width=\linewidth]{box_fps.pdf}
        \caption{Frames per second per ablation step.}
        \label{fig:box_fps}
    \end{subfigure}
    \caption{Per-run distribution of load time and FPS across all scenarios and styles, by ablation step.
    Load-time variance compresses sharply at \texttt{tile\_shave}, while the FPS spread widens past the same step as per-style headroom diverges.}
    \label{fig:box}
\end{figure}

\begin{figure}[tb]
    \centering
    \includegraphics[width=\linewidth]{waterfall_gzip_bytes.pdf}
    \caption{Cumulative ablation of mean style gzip bytes across all 19 steps.
    Most of the wire-size win lands by metadata stripping.
    The \texttt{tile\_shave\_only} spike is the unoptimized-style comparison configuration, which the subsequent steps recover and surpass.}
    \label{fig:waterfall}
\end{figure}

\subsection{Cross-style generalization}\label{sec:results:generalization}
Applied with static passes only to the \num{13} unseen Maputnik styles, the optimizer delivered a median raw style-size reduction of \qty{31.4}{\percent} (range \qtyrange{3.2}{67.7}{\percent}), narrowing to a median of \qty{11.4}{\percent} under gzip (range \qtyrange{3.9}{22.4}{\percent}), with no style regressing (Figure~\ref{fig:cross_style}).
The top reducers were styles with verbose filter logic or repetitive property structures (LocationIQ Streets at \qty{67.7}{\percent}, Americana at \qty{59.2}{\percent}, MapTiler Toner at \qty{57.5}{\percent}), while already-minified basemaps such as Stadia Outdoors (\qty{3.2}{\percent}) and the AWS pair (\qty{10.6}{\percent}) left little to remove.
Initial style complexity did not reliably predict achievable reduction (Pearson $r = 0.35$, $n = 13$, $p = 0.24$), a relationship dominated by a single high-leverage point such that dropping Americana flips the sign to $r = -0.54$.
Because the generalization corpus carries no tile data, stats-driven and data-level passes were not applied, so these figures are a lower bound on the full potential.

\begin{figure}[tb]
    \centering
    \includegraphics[width=\linewidth]{cross_style_reduction.pdf}
    \caption{Per-style raw, gzip, and brotli style-size reduction across the \num{13} unseen Maputnik styles, sorted by raw reduction.
    Raw reductions range from \qty{3.2}{\percent} to \qty{67.7}{\percent}, with no style regressing.}
    \label{fig:cross_style}
\end{figure}

\subsection{Threats to validity}\label{sec:results:threats}
All rendering benchmarks ran on a single hardware and browser configuration with ANGLE/Vulkan, so results may differ on other hardware or graphics backends, and the 15-run median may not capture rare events such as thermal throttling.
The evaluation is limited to MapLibre GL JS and the OMT schema, and the 13 benchmark styles, though diverse, may not cover all authoring patterns.
Network conditions were simulated through a fixed bandwidth throttle and a caching proxy rather than real content delivery network latency, which also masks real-world cache behavior.
The cumulative ablation design means each pass's reported marginal contribution depends on the order of preceding passes, and tail metrics (p99 frame time, jank count) carry higher variance than central-tendency metrics.
Because the \num{18} scenarios share hardware, tile data, and styles, they are not fully independent, so we read the reported $p$-values as descriptive evidence of large effect sizes rather than strict inferential tests, and apply no multiple-comparisons correction across the full ablation grid.
Visual-equivalence checks across the deep corpus, all scenarios, and zoom levels 0 to 14 confirmed every pass stays within the per-test pixel-difference tolerance of Section~\ref{sec:methods:problem}.
